\subsection{Frontier Infrastructure: Working Where the Verifier Is Unforgiving}
\label{sec:infra-vertical}

The campaigns above are evaluated by a proof checker, a venue format, and an
official benchmark harness. This subsection covers the vertical where we have the
most external evidence and where an unsupported claim has the shortest life: GPU
kernel and serving-infrastructure engineering. A kernel is either bit-accurate
against a reference or it is not; it is either faster on the same hardware with the
same shapes or it is not---the pairing of correctness and speedup that
KernelBench formalizes as an evaluation target~\citep{ouyang2025kernelbench}; and when the work is submitted upstream, the party
deciding whether it is acceptable is a maintainer with no stake in this report.

The vertical spans both ends of the hardware range: datacenter accelerators where
the constraint is throughput, and consumer devices where the constraint is whether
the model fits at all.

This domain is also the sharpest test of the Driver framing, because it is where the
objective moves the most. None of the campaigns below could state their target
precisely at the start. The FlashAttention-4 adaptation did not know which of six
mask families would fit the tile scheduler, and spent its first several rounds
discovering that its initial data path was catastrophically wrong. The SGLang
integration did not know that a cross-batch determinism defect existed until
batch size eight produced drifting images. The correctness fix described below began
as a performance investigation and became a correctness one. Each is a
dense-intelligence task in the sense of Section~\ref{sec:problem}: invention rather
than retrieval, many dependent iterations, and a task-native verifier strong enough
that the runtime can judge itself without appeal to its own narrative.

\subsubsection{Diffusion-LLM Masks on FlashAttention-4}

Diffusion language models do not use causal attention. They require block-causal,
prefix-full, prefix-causal, prefix-hole, sliding-window, and cache-only mask
families, plus compositions such as prefix-full with holes under a sliding window.
The task was to carry all of them into a current-generation FlashAttention-4 CuTe
DSL kernel across both Ampere SM80 and Hopper SM90 paths, with paged and dense KV
layouts, without surrendering the performance that motivated using FA4 in the first
place.

The result is \FlashDAMissions{} missions and \FlashDACalls{} calls over
\FlashDACompute{} hours of model compute, producing 36 changed files
(\code{+5,814/-14}). All 19 parity cases pass across both architectures, paged and
dense dispatch, composed masks, and cache-only inputs, with maximum deviation of
$3.9\times10^{-3}$ at bf16 rounding scale and cosine similarity at or above
$0.9999994$; CUDA Graph replay is bit-identical to direct invocation. Against the
production Triton path the adaptation is $2.5$--$4.2\times$ faster end to end, or
$1.6$--$2.6\times$ when both sides enable CUDA Graph. Against upstream native FA4 it
is slower by only \FlashDAGapLow--\FlashDAGapHigh\%, and that residual is almost
entirely a fixed 3.6--5\,\textmu s Python wrapper cost: the attention kernel itself
runs within 2--4\,\textmu s of native.

Two facts about the process matter more than the endpoint. First, the campaign's
early dense-plus-block-sparsity route was not slightly worse but $4.9$--$29.6\times$
slower than native, with dense KV gather alone reaching 884.5\,\textmu s at decode.
The final result came from recognising that this data path was wrong and abandoning
it, not from optimising within it---and the rejected route, with its evidence, is
retained in the skill library rather than discarded. Second, the entire campaign
cost under \FlashDACostCNY{} CNY of model spend, distributed across three vendors by
role: an open-weight model carried 82\% of Engineer compute while separate models
served Manager triage and Planner decomposition, all under one role contract. A
specialist infrastructure team had previously worked this direction for over three
months; that figure comes from a practitioner interview covering work of different
scope, not a matched trial.

\subsubsection{A Unified Model into a Production Serving Stack}

Integrating SenseNova U1 into SGLang required five workstreams that are usually
staffed separately: architecture and weight mapping for language and multimodal
understanding; native bounded flow-matching image generation; scheduler-owned
text-to-image-to-text interleaving with correct KV lifecycle; CUDA Graph
flow-prefill with concurrency determinism; and an OpenAI-compatible API with
cancellation, disconnect, and multipart limits.

The delivered change is 36 files and \code{+11,263/-72} lines across 14 commits.
Acceptance was established on real hardware rather than by structural review: 1,116
tensors load with zero missing or unknown, visual question answering is exact on
160 of 160 generated tokens, concurrency-8 is exact on 8 of 8, and throughput
improves $5.108\times$ at batch size 8 with $2.276\times$ on text-to-image and
$3.618\times$ on image-conditioned generation. A cross-batch determinism defect that
caused image drift at batch size 8 was found, localised, and fixed, with the
invariant then holding across all batch sizes.

The comparison in Section~\ref{sec:delivery} is against this task: one engineer
working with a turn-by-turn coding agent invested over \SGLangHumanHours{} hours
without completing it, while a blind agent run stopped after roughly 1 hour 21
minutes with a CPU-oriented draft carrying no real weights, no GPU parity, and no
benchmarks. \systemname{} completed it within a \SGLangArgusHours{}-hour envelope
under an operator who was not an infrastructure specialist. What separates the
outcomes is not the ability to write plausible integration code but the willingness
to keep going through the expensive second half: real checkpoint loading, exact
parity, concurrency determinism, service-level safety, and the revisions that
external review demands.

\subsubsection{Upstream Review by Maintainers Who Owe Us Nothing}

Kernels authored by \systemname{} were submitted to
\code{flash-linear-attention}, a widely used implementation library for
linear-attention architectures. Two properties make this the strongest external
evidence in this report: the acceptance standard is written by someone else---the
repository requires dependent-test discovery and same-hardware before-and-after
benchmarks---and the reviewers had no involvement in producing the work.

Four submissions are recorded here, spanning performance work and correctness
repair. A TileLang RWKV6 forward-intra backend (PR~\#1045) improved forward latency from 0.199\,ms to
0.168\,ms ($1.18\times$) and forward-plus-backward from 0.900\,ms to 0.747\,ms
($1.21\times$) on an H100 NVL at $B{=}8$, $T{=}1024$, $H{=}8$, $D{=}64$, gated by 14
backend tests and a 13-case frozen verification gate, with unsupported dtypes,
dimensions, and layouts rejected by a backend verifier rather than silently
mishandled. Its submission states plainly that the implementation, optimisation
loop, correctness validation, and performance evidence were completed autonomously
by \systemname{}, with only the target and hardware constraint chosen by the
operator. That sentence was reviewed and accepted along with the code by an outside
maintainer, which is a different kind of claim from one asserted in a report about
oneself. It received no inline change requests.

A second submission (PR~\#1128) accelerates \code{chunk\_kda} training with four TileLang
kernels covering intra forward, $dAv$, fused $dq/kg$, and intra backward, reaching
$1.29\times$ over Triton across the four stages on B200. The instructive detail is
what was not claimed: the best single stage reached $1.541\times$, but automatic
dispatch was enabled only for the one workload with both verified correctness and a
repeatable end-to-end gain, which measures $1.078$--$1.099\times$. The larger number
was available and was not used to characterise the result.

The third (PR~\#1109) is two lines. On SM100, \code{chunk\_kda} backward autotuning triggered an
illegal memory access that poisoned the CUDA context, so the test file could not run
to completion; excluding \code{BK==32} configurations at 4 and 8 warps restores a
full pass at 76 tests. There is no speedup to report, because before the fix nothing
could be measured at all. The reviewing maintainer independently re-derived that 24
autotune candidates survive the filter, confirmed no empty-configuration risk,
verified that the described scope matched the diff, and diagnosed two red CI checks
as infrastructure flakes rather than test failures, then approved with two
non-blocking suggestions and no required changes. The fix had also been narrowed by
its author from all Blackwell devices to the measured SM100 target---the same
discipline of publishing the perimeter of a claim that appears in the ACE-2
certificate of Section~\ref{sec:silicon-vertical}.

A fourth submission (PR~\#1114) parallelizes the long-sequence \code{AttnRes}
reduction, specializing forward and backward kernels on the exact residual count and
distributing the $dq$, $dw$, and optional $dow$ reductions across the flattened token
dimension for $N\geq16384$. Across five bfloat16 shapes on B200 it reaches a
geometric-mean speedup of $1.102\times$, best case $1.237\times$. The submission
reports its worst row, $1.033\times$, alongside the mean, and states that no broader
hardware coverage or end-to-end model speedup is claimed.

\subsubsection{The Same Discipline at the Other End of the Hardware Range}
\label{sec:local-deployment}

The campaigns above optimize kernels on datacenter accelerators. This one inverts
the constraint: the model is fixed, and the hardware is what an individual actually
owns. MiniMax-H3 is an audio-video generation
model whose BF16 diffusion transformer alone occupies roughly 62\,GiB, well beyond
the memory of the two targets below. The task was not to quantize or distill it into
something smaller, but to make the original weights run.

\paragraph{Apple M4 Pro, 24\,GB unified memory.}
\systemname{} adapted the upstream inference path to Apple MLX with block-wise
streaming load and prompt release, so the full BF16 diffusion transformer and BF16
text encoder are never resident at once. The system produces
$1344\times768$ video, 124 frames at 24\,FPS with 5.17 seconds of stereo audio, in
47\,minutes 58.7\,seconds end to end at a peak memory footprint of approximately
15.8\,GB---roughly a quarter of the model's nominal size, on a laptop-class device.

\paragraph{One RTX A6000, 48\,GB.}
The same model was brought to a single Ampere-class workstation GPU at full FL2VA
checkpoint fidelity, $1344\times768$, 124 frames, 24\,FPS, 32\,kHz stereo. The
campaign then optimized it under measurement discipline. A BF16 dense warm baseline
over $N{=}10$ runs establishes 1{,}792.202\,s. An eight-step Turbo route reaches
290.998\,s, a $6.159\times$ speedup, and is adopted as the practical default. A
sparse-attention route at $r{=}8$ is reported separately at a formal $N{=}10$
HTTP-time improvement of 15.203\%, and a matched formal $N{=}10$ comparison on
30-second final audio-video yields $+4.326\%$.

\paragraph{What makes this a runtime result rather than a port.}
Both targets are public repositories, and both record what was rejected as well as
what was kept. Candidates that degraded output quality, produced no gain on the
critical path, or rested on $N{=}1$ evidence are marked rejected rather than folded
into the headline number, and quality failures are published alongside the speedups.
The distinction matters because the two accepted routes are not equivalent: the
Turbo path is an approximate decoding schedule and is labeled as such, not presented
as lossless BF16 acceleration. Acceptance ran through paired $N{=}10$ comparisons on
the same physical GPU, a 24-case quality suite, and memory and audio-video integrity
checks.

This campaign exercises a different part of the runtime than kernel optimization.
There is no leaderboard and no official verifier; it had to construct its own
measurement protocol, decide which improvements were real, and publish the boundary
of each claim. That is the same discipline the ACE-2 certificate applies to hardware
and the materials campaign applies to chemistry, in a setting where the deliverable
is simply that a model too large for the device runs on it.

\subsubsection{Why This Vertical Is Evidence About the Runtime}

These campaigns share a shape that a score-climbing system does not produce. One
result is a rejected route retained as reusable knowledge. One is a deliberate
choice to report the smaller honest number over the larger available one. One is a
correctness fix with no performance story, whose entire value is that a test suite
that could not finish now finishes. Together with three merged or reviewed upstream
submissions and a serving integration validated on real hardware, they indicate that
the unattended hours produce work that survives contact with people who did not
write it---which is the property that distinguishes delegated autonomy from
generated volume.
